\chapter{Finite Fields}

In Chapter~5 we saw that a field is a place where you can add, subtract, multiply, and divide, and where the usual algebra you learned in school keeps working: zero times anything is zero, $a(b+c)=ab+ac$, and every nonzero element has a reciprocal. The rational numbers $\Q$, the reals $\R$, and the complex numbers $\C$ are all fields. They are also all infinite, and a computer cannot store an arbitrary real number. Cryptography runs on machines that manipulate fixed-size chunks of bits, so we need fields with only finitely many elements: \emph{finite fields}.

This chapter builds the two kinds of finite field that nearly all of modern
cryptography stands on. The first, $\F_p = \Zn{p}$ for a prime $p$, you have
already met --- it is arithmetic modulo a prime, and it is
where Diffie--Hellman, RSA's exponent group, and elliptic curves live. The
second, $\F_{2^n}$, is stranger and more useful: a field whose elements are
bit-strings, whose addition is XOR, and whose multiplication mixes bits in a
way that is cheap in hardware and algebraically rich. AES lives in
$\F_{2^8}$, and by the end of this chapter you will have built that exact
field by hand, multiplied two bytes in it, and understood why the AES
designers chose it.

The promise of the chapter: every finite field has a prime-power number of
elements $q = p^n$, there is essentially only one field of each size, and we
can construct it concretely with nothing more than polynomials and long
division. Once you see that, $\F_{2^8}$ stops looking like an opaque recipe in
the AES specification and becomes something you can compute with on paper.

\chapterline

\section{What a Field Is, Quickly}
\label{ch06:what-is-field}

We will lean on the previous chapter's definitions, but let us restate the
one that matters so the chapter stands on its own.

\begin{definition}[Field]
\label{ch06:def-field}
A \emph{field} is a set $\F$ with two operations, addition $+$ and
multiplication $\cdot$, such that:
\begin{itemize}
  \item $(\F, +)$ is an abelian group with identity $0$ (you can add and
        subtract; addition commutes);
  \item $(\F \setminus \set{0}, \cdot)$ is an abelian group with identity
        $1$ (every nonzero element has a multiplicative inverse; you can
        multiply and divide);
  \item multiplication distributes over addition: $a\cdot(b+c) = a\cdot b +
        a\cdot c$.
\end{itemize}
\end{definition}

The single clause that does all the heavy lifting is ``every nonzero element
has an inverse.'' That is what separates a field from a mere ring like
$\Zn{n}$ for composite $n$. In $\Zn{6}$ the element $2$ has no inverse: there
is no $y$ with $2y \equiv 1 \pmod 6$, because $2$ and $6$ share the factor
$2$. So $\Zn{6}$ is a ring but not a field. The moment the modulus is prime,
\emph{every} nonzero element is coprime to it, every nonzero element gets an
inverse, and we have a field.

\begin{intuitionbox}
A field is the smallest setting in which ``solve $ay = b$ for $y$'' always
has a unique answer when $a \neq 0$: just multiply both sides by $a^{-1}$.
Without inverses, no division, and almost every cryptographic construction
that needs to ``undo'' an operation would fall apart. AES decryption, Shamir
secret sharing, Reed--Solomon codes, and polynomial commitment schemes all
rely on being able to invert and divide inside a finite field.
\end{intuitionbox}

\subsection{The first family: $\F_p$}

\begin{theorem}[Prime moduli give fields]
\label{ch06:thm-fp}
For a prime $p$, the ring $\Zn{p} = \set{0,1,\dots,p-1}$ with addition and
multiplication mod $p$ is a field. We write it $\F_p$ (or $\mathrm{GF}(p)$,
``Galois field of order $p$'').
\end{theorem}

\begin{proof}
$\Zn{p}$ is already a commutative ring with $1$. We only need multiplicative
inverses for nonzero elements. Take $a \in \set{1,\dots,p-1}$. Since $p$ is
prime and $0 < a < p$, we have $\gcd(a,p)=1$. By B\'ezout's identity (the
extended Euclidean algorithm from Chapter~4) there are
integers $u,v$ with $ua + vp = 1$. Reducing mod $p$ gives $ua \equiv 1
\pmod p$, so $u \bmod p$ is the inverse of $a$. Every nonzero element is
invertible, so $\F_p$ is a field.
\end{proof}

\begin{example}[Inverses in $\F_7$]
\label{ch06:ex-f7}
In $\F_7$, what is $3^{-1}$? We want $3y \equiv 1 \pmod 7$. Trying
$y=5$: $3\cdot 5 = 15 = 14 + 1 \equiv 1$. So $3^{-1} = 5$. The full inverse
table for $\F_7$:
\[
1^{-1}=1,\quad 2^{-1}=4,\quad 3^{-1}=5,\quad 4^{-1}=2,\quad 5^{-1}=3,\quad
6^{-1}=6.
\]
Notice $6 = -1$ is its own inverse, as is $1$. Everything else pairs up.
Division is now legal: $\tfrac{2}{3}$ in $\F_7$ means $2 \cdot 3^{-1} = 2
\cdot 5 = 10 \equiv 3$.
\end{example}

\begin{connectionbox}
Computing $a^{-1}$ in $\F_p$ is exactly the extended Euclidean algorithm,
which runs in $O(\log p)$ steps --- see the algorithms companion for the cost
analysis. An alternative uses Fermat's little theorem: $a^{p-1}\equiv 1$, so
$a^{-1} \equiv a^{p-2} \pmod p$, computed by fast modular exponentiation in
$O(\log p)$ multiplications. The Euclidean route is usually faster; the
Fermat route is constant-time-friendlier, which matters for side-channel
resistance (see Chapter~37).
\end{connectionbox}

\section{Why Prime Powers, and Nothing Else}
\label{ch06:prime-powers}

So $\F_p$ exists for every prime. Can we have a field with, say, $6$
elements? Or $4$? The answer is one of the cleanest theorems in algebra.

\begin{theorem}[Classification of finite fields]
\label{ch06:thm-classification}
A finite field has exactly $q = p^n$ elements for some prime $p$ and integer
$n \geq 1$. Conversely, for every prime power $p^n$ there is a finite field
with that many elements, and any two finite fields with the same number of
elements are isomorphic --- structurally identical up to relabeling their elements.
\end{theorem}

We will not prove the whole theorem, but the ``why $p^n$'' half is short and
illuminating, so here it is.

\begin{whybox}
Let $\F$ be a finite field. Look at the elements you get by repeatedly adding
$1$ to itself: $1, 1+1, 1+1+1, \dots$ Because $\F$ is finite, this sequence
must eventually repeat, and the first time you return to $0$ defines a number
$p$ with $\underbrace{1+1+\dots+1}_{p}=0$. This $p$ is called the
\emph{characteristic}. It must be prime: if $p = ab$ with $1<a,b<p$, then
$(\underbrace{1+\dots+1}_{a})(\underbrace{1+\dots+1}_{b}) = 0$, a product of
two nonzero elements equalling zero --- impossible in a field, where nonzero
times nonzero is nonzero (otherwise one of them would be a zero-divisor and
could not have an inverse).

Now the copy of $\set{0,1,\dots,p-1}$ sitting inside $\F$ is a sub-field
isomorphic to $\F_p$. The whole field $\F$ is a vector space over this
$\F_p$: you can add field elements and scale them by elements of $\F_p$, and
all the vector-space axioms hold. A finite-dimensional vector space over
$\F_p$ of dimension $n$ has exactly $p^n$ elements --- choose each of the $n$
coordinates from $p$ options. Hence $\abs{\F} = p^n$.
\end{whybox}

That argument also tells us \emph{how} to build $\F_{p^n}$: realize it as an
$n$-dimensional vector space over $\F_p$ with a compatible multiplication.
The trick for the multiplication is polynomials.

\begin{intuitionbox}
Why can't there be a field with $6$ elements? Because $6 = 2\cdot 3$ is not a
prime power. There is a field of size $4=2^2$, a field of size $8=2^3$, a
field of size $9 = 3^2$, but none of size $6$, $10$, $12$, or $15$. The gaps
are exactly the non-prime-powers. This is not a limitation of cleverness; it
is a theorem.
\end{intuitionbox}

\section{Polynomials over a Field}
\label{ch06:polynomials}

To build $\F_{p^n}$ we copy the construction of $\F_p$ itself, one level up.
Recall how $\F_p$ was made: take the integers $\Z$, pick a prime $p$, and
work modulo $p$. The primality of $p$ was what made the quotient a field.

We will do the same with \emph{polynomials} in place of integers. Fix a base
field $\F_p$ and consider polynomials in one variable $x$ with coefficients
in $\F_p$:
\[
\F_p[x] = \set{\, a_{d}x^{d} + \dots + a_1 x + a_0 : a_i \in \F_p \,}.
\]
You add them coefficient-by-coefficient and multiply them the usual way,
reducing coefficients mod $p$. This set $\F_p[x]$ behaves astonishingly like
$\Z$:

\begin{itemize}
  \item There is a notion of \emph{degree} (the highest power with nonzero
        coefficient), playing the role that size/magnitude plays for
        integers.
  \item There is \emph{division with remainder}: given $f$ and a nonzero
        $g$, there are unique $q$ (quotient) and $r$ (remainder) with $f =
        qg + r$ and $\deg r < \deg g$. This is just polynomial long division.
  \item There are \emph{irreducible} polynomials --- ones that do not factor
        into lower-degree pieces over $\F_p$. These are the ``primes'' of
        $\F_p[x]$.
\end{itemize}

\begin{intuitionbox}
The dictionary to keep in your head:
\[
\Z \leftrightarrow \F_p[x], \qquad
\text{prime} \leftrightarrow \text{irreducible polynomial}, \qquad
\Zn{p} \leftrightarrow \F_p[x]/(m(x)).
\]
Working modulo a prime integer gave a field; working modulo an irreducible
polynomial will give a field too. The degree of the polynomial controls how
big the field is.
\end{intuitionbox}

\subsection{Irreducibility with small numbers}

Over $\F_2$, the coefficients are just $0$ and $1$, so a polynomial is
nothing more than a bit-string telling you which powers of $x$ are present.
For example $x^3 + x + 1$ is the bit-string $1011$.

A polynomial of degree $\geq 1$ is reducible over $\F_2$ if it factors. The
quickest first test: a polynomial has the factor $x$ iff its constant term is
$0$, and it has the factor $(x+1)$ iff it has an even number of terms (iff
$x=1$ is a root, since over $\F_2$ a root at $1$ means the coefficients sum to
$0$).

\begin{example}[Cataloguing low-degree irreducibles over $\F_2$]
\label{ch06:ex-irreducibles}
\emph{Degree 1:} $x$ and $x+1$ are both irreducible (degree-1 polynomials
always are).

\emph{Degree 2:} the candidates are $x^2, x^2+1, x^2+x, x^2+x+1$. The first
three have a root in $\F_2$ ($x^2$ and $x^2+x$ have root $0$; $x^2+1 =
(x+1)^2$ has root $1$). Only $x^2+x+1$ has no root in $\F_2$ --- check:
plug in $0$ gives $1$, plug in $1$ gives $1+1+1=1$. So $x^2+x+1$ is the
unique irreducible quadratic over $\F_2$, and it builds $\F_4$.

\emph{Degree 3:} a cubic over $\F_2$ is reducible iff it has a root (a cubic
that factors must have a linear factor). Testing the $8$ cubics, exactly two
have no root: $x^3+x+1$ and $x^3+x^2+1$. Either one builds $\F_8$.
\end{example}

\begin{walkbox}
Is $x^4 + x + 1$ irreducible over $\F_2$? It has no root: $f(0)=1$,
$f(1)=1+1+1=1$. So it has no linear factor. But a degree-4 polynomial could
still factor as two irreducible quadratics. The only irreducible quadratic
over $\F_2$ is $x^2+x+1$, and $(x^2+x+1)^2 = x^4 + x^2 + 1 \neq x^4+x+1$.
No linear factor, not the square of the lone irreducible quadratic --- so
$x^4+x+1$ is irreducible. (This polynomial builds $\F_{16}$ and shows up in
CRC checksums.)
\end{walkbox}

\section{Building $\F_4$ by Hand}
\label{ch06:build-f4}

Let us construct the smallest non-prime finite field completely, because
once you have seen $\F_4$ in full, $\F_{2^8}$ is just the same idea with
bigger bit-strings.

Take $\F_2$ and the irreducible polynomial $m(x) = x^2 + x + 1$. Define
\[
\F_4 = \F_2[x] / (m(x)),
\]
the polynomials over $\F_2$ taken modulo $m(x)$. Because $m$ has degree $2$,
every remainder has degree $< 2$, so the elements are
\[
\F_4 = \set{\, 0,\ 1,\ x,\ x+1 \,},
\]
exactly four of them --- $2^2$, as promised. Often $x$ is written $\alpha$
and people call $\alpha$ ``a root of $m$.''

Addition is XOR of the coefficient vectors (since we are over $\F_2$):
\[
(x) + (x+1) = 1, \qquad (x+1)+(x+1) = 0, \qquad 1 + 1 = 0.
\]

Multiplication multiplies the polynomials and then reduces mod $m(x)$. The
key relation comes from $m(x)=0$ in the quotient:
\[
x^2 + x + 1 = 0 \quad\Longrightarrow\quad x^2 = x + 1
\]
(remember $-1 = +1$ over $\F_2$). So whenever an $x^2$ appears, replace it
with $x+1$. Watch:
\[
x \cdot x = x^2 = x + 1,
\]
\[
x \cdot (x+1) = x^2 + x = (x+1) + x = 1,
\]
\[
(x+1)(x+1) = x^2 + 1 = (x+1) + 1 = x.
\]

\begin{center}
\begin{tabular}{c|cccc}
$\cdot$ & $0$ & $1$ & $x$ & $x{+}1$ \\
\hline
$0$     & $0$ & $0$   & $0$   & $0$ \\
$1$     & $0$ & $1$   & $x$   & $x{+}1$ \\
$x$     & $0$ & $x$   & $x{+}1$ & $1$ \\
$x{+}1$ & $0$ & $x{+}1$ & $1$ & $x$ \\
\end{tabular}
\end{center}

Read off the inverses from the table: $x^{-1} = x+1$ (their product is $1$),
and $(x+1)^{-1} = x$. Every nonzero element has an inverse, so this is a
genuine field --- and it has $4$ elements, the field you cannot build out of
$\Zn{4}$ (which has the zero-divisor $2\cdot 2 = 0$).

\begin{trapbox}
$\F_4$ is \emph{not} $\Zn{4}$. In $\Zn 4$, $1+1 = 2$ and $2+2 = 0$; the
additive group is cyclic of order $4$. In $\F_4$, every element added to
itself is $0$ (characteristic $2$); the additive group is $\Z_2 \times
\Z_2$, not $\Z_4$. They have the same size and are both rings, but only
$\F_4$ is a field. Mixing them up is one of the most common beginner errors
with $\mathrm{GF}(2^n)$.
\end{trapbox}

\begin{figure}[H]
\centering
\begin{tikzpicture}[font=\small,>=stealth,scale=1.1]
  % Z_4 structure
  \begin{scope}[shift={(0,0)}]
    \node[font=\bfseries,sourceblue!90!black] at (0, 1.8) {Ring $\Zn{4}$};
    \node[font=\footnotesize,gray] at (0, 1.4) {Additive Group: Cyclic $C_4$};
    
    \node[draw=sourceblue!80,fill=sourceblue!10,circle,thick,minimum size=0.8cm,inner sep=0pt] (z0) at (0, 0.7) {$0$};
    \node[draw=sourceblue!80,fill=sourceblue!10,circle,thick,minimum size=0.8cm,inner sep=0pt] (z1) at (0.7, 0) {$1$};
    \node[draw=sourceblue!80,fill=sourceblue!10,circle,thick,minimum size=0.8cm,inner sep=0pt] (z2) at (0, -0.7) {$2$};
    \node[draw=sourceblue!80,fill=sourceblue!10,circle,thick,minimum size=0.8cm,inner sep=0pt] (z3) at (-0.7, 0) {$3$};

    \draw[->,thick,sourceblue!80] (z0) to[bend left=20] node[above right,font=\tiny]{$+1$} (z1);
    \draw[->,thick,sourceblue!80] (z1) to[bend left=20] node[below right,font=\tiny]{$+1$} (z2);
    \draw[->,thick,sourceblue!80] (z2) to[bend left=20] node[below left,font=\tiny]{$+1$} (z3);
    \draw[->,thick,sourceblue!80] (z3) to[bend left=20] node[above left,font=\tiny]{$+1$} (z0);
  \end{scope}

  % Divider
  \draw[dashed,gray!50] (2.5, 1.8) -- (2.5, -1.2);

  % F_4 structure
  \begin{scope}[shift={(5,0)}]
    \node[font=\bfseries,targetgreen!90!black] at (0, 1.8) {Field $\F_4$};
    \node[font=\footnotesize,gray] at (0, 1.4) {Additive Group: Hypercube $\Z_2 \times \Z_2$};
    
    \node[draw=targetgreen!80,fill=targetgreen!10,thick,rounded corners=2pt,minimum width=0.9cm,minimum height=0.6cm] (f0) at (-1, 0.6) {$0$};
    \node[draw=targetgreen!80,fill=targetgreen!10,thick,rounded corners=2pt,minimum width=0.9cm,minimum height=0.6cm] (f1) at (1, 0.6) {$1$};
    \node[draw=targetgreen!80,fill=targetgreen!10,thick,rounded corners=2pt,minimum width=0.9cm,minimum height=0.6cm] (fx) at (-1, -0.6) {$x$};
    \node[draw=targetgreen!80,fill=targetgreen!10,thick,rounded corners=2pt,minimum width=0.9cm,minimum height=0.6cm] (fx1) at (1, -0.6) {$x{+}1$};

    \draw[<->,thick,targetgreen!80] (f0) -- node[above,font=\tiny]{$+1$} (f1);
    \draw[<->,thick,targetgreen!80] (fx) -- node[below,font=\tiny]{$+1$} (fx1);
    \draw[<->,thick,targetgreen!80] (f0) -- node[left,font=\tiny]{$+x$} (fx);
    \draw[<->,thick,targetgreen!80] (f1) -- node[right,font=\tiny]{$+x$} (fx1);
  \end{scope}
\end{tikzpicture}
\caption{Additive structure of $\Zn{4}$ vs $\F_4$. The integer ring $\Zn{4}$ forms a single cycle of 4 elements under addition, meaning $1+1 = 2 \neq 0$. The Galois field $\F_4$ forms a 2D hypercube of degree-1 polynomials over $\F_2$; every element is its own inverse, meaning $x+x = 0$ and $1+1 = 0$, giving it a fundamentally different algebraic topology.}
\label{ch06:fig-f4-vs-z4}
\end{figure}

\begin{historybox}
These fields are named for \'Evariste Galois (1811--1832), who published ``Sur la th\'eorie des nombres'' in 1830 introducing finite fields of prime-power order $p^n$ by adjoining roots of irreducible polynomials modulo $p$. In 1893, the American mathematician Eliakim Hastings Moore proved that all finite fields of the same order $p^n$ are isomorphic and coined the term \emph{Galois field}. The notation $\mathrm{GF}(q)$ is standard in engineering, while mathematicians write $\F_q$: for the same object, AES documents say $\mathrm{GF}(2^8)$ and algebra texts say $\F_{256}$. In modern computing, $\mathrm{GF}(2^8)$ powers the AES S-box and $\mathrm{GF}(2^{128})$ powers the GHASH authenticator in AES-GCM, though quasi-polynomial discrete-logarithm algorithms in characteristic~$2$ have pushed public-key cryptography away from binary fields toward prime-order fields.
\end{historybox}

\section{The Field AES Uses: $\F_{2^8}$}
\label{ch06:f256}

A byte is eight bits, and the natural arena for byte-oriented cryptography is
a field of $256 = 2^8$ elements. We identify a byte $b_7 b_6 \cdots b_1 b_0$
with the polynomial
\[
b_7 x^7 + b_6 x^6 + \dots + b_1 x + b_0 \in \F_2[x],
\]
a polynomial of degree at most $7$. For instance the byte
\texttt{0x57} $= 0101\,0111$ corresponds to $x^6 + x^4 + x^2 + x + 1$.

To make these $256$ bit-strings into a field, we pick an irreducible
polynomial of degree $8$ over $\F_2$ and reduce modulo it. AES fixes the
choice
\[
\boxed{\ m(x) = x^8 + x^4 + x^3 + x + 1\ }
\qquad(\texttt{0x11B} \text{ in hex}),
\]
sometimes called the Rijndael polynomial. So
\[
\F_{2^8} = \F_2[x]/(m(x)),
\]
with elements the $256$ bytes, addition $=$ XOR, and multiplication $=$
polynomial multiplication followed by reduction mod $m(x)$.

\begin{whybox}
Why is $x^8+x^4+x^3+x+1$ irreducible, and why this one? Irreducibility you
can verify by trial division: it has no root in $\F_2$ ($f(0)=1$ and
$f(1)=1$), and you check it shares no factor with the irreducible polynomials
of degree $2,3,4$. The AES designers needed \emph{some} degree-8 irreducible;
there are $30$ of them. They picked one with few terms so that reduction is
cheap: a smaller polynomial means fewer XORs when you reduce. The field
$\F_{2^8}$ you get is the same regardless of which irreducible you choose (all
fields of size $256$ are isomorphic by \autoref{ch06:thm-classification});
only the bit-level recipe for multiplication changes.
\end{whybox}

\begin{figure}[H]
\centering
\begin{tikzpicture}[font=\small,>=stealth,scale=0.92]
  % Panel 1: Byte-to-Polynomial mapping
  \begin{scope}[shift={(0,3.2)}]
    \node[font=\bfseries\footnotesize,sourceblue,anchor=east] at (-0.6,0) {Byte / Polynomial:};
    \foreach \i/\bit/\term in {7/0/x^7, 6/1/x^6, 5/0/x^5, 4/1/x^4, 3/0/x^3, 2/1/x^2, 1/1/x, 0/1/1} {
      \pgfmathsetmacro\xpos{(7-\i)*1.1}
      \ifnum\bit=1
        \node[draw=sourceblue!80,fill=sourceblue!20,thick,rounded corners=2pt,minimum width=0.95cm,minimum height=0.6cm,font=\bfseries] at (\xpos,0) {\bit};
        \node[below,font=\tiny\bfseries,sourceblue!90!black] at (\xpos,-0.35) {$\term$};
      \else
        \node[draw=gray!50,fill=gray!10,rounded corners=2pt,minimum width=0.95cm,minimum height=0.6cm,font=\bfseries,text=gray] at (\xpos,0) {\bit};
        \node[below,font=\tiny,gray] at (\xpos,-0.35) {$0$};
      \fi
      \node[above,font=\tiny,gray] at (\xpos,0.35) {$b_\i$};
    }
    \node[right,font=\footnotesize\bfseries,sourceblue!90!black] at (8.0,0) {$= \texttt{0x57} \in \F_{2^8}$};
  \end{scope}

  % Divider
  \draw[dashed,gray!40] (-2.2,2.2) -- (11.0,2.2);

  % Panel 2: Carry-less XOR addition
  \begin{scope}[shift={(0,1.2)}]
    \node[font=\bfseries\footnotesize,targetgreen] at (-1.5,0) {Addition (XOR):};
    \node[font=\scriptsize\ttfamily] at (1.5,0.4) {A = 0101 0011 (0x53)};
    \node[font=\scriptsize\ttfamily] at (1.5,-0.1) {B = 1100 1010 (0xCA)};
    \draw[thick,gray!60] (0.2,-0.35) -- (2.8,-0.35);
    \node[font=\scriptsize\ttfamily\bfseries,targetgreen!90!black] at (1.5,-0.65) {C = 1001 1001 (0x99)};
    \node[font=\tiny,align=left] at (6.5,-0.1) {
      $\bullet$ Bitwise XOR on coefficients mod 2.\\
      $\bullet$ \textbf{Zero carry propagation:} isolated bit-level logic.
    };
  \end{scope}

  % Divider
  \draw[dashed,gray!40] (-2.2,0.1) -- (11.0,0.1);

  % Panel 3: Modular Reduction by Rijndael polynomial
  \begin{scope}[shift={(0,-1.5)}]
    \node[font=\bfseries\footnotesize,quotientorange,anchor=east] at (0.2,0) {Modular Reduction:};
    \node[draw=quotientorange!80,fill=quotientorange!10,rounded corners=3pt,thick,align=left,font=\tiny,inner sep=4pt] at (4.0,0) {
      \textbf{Irreducible Polynomial (AES modulus):}\\
      $m(x) = x^8 + x^4 + x^3 + x + 1 \equiv 0$\\
      $\implies \mathbf{x^8 \equiv x^4 + x^3 + x + 1} \quad (\texttt{0x1B})$
    };

    \node[draw=structurepurple!80,fill=structurepurple!10,rounded corners=3pt,thick,align=left,font=\tiny,inner sep=4pt] at (8.6,0) {
      \textbf{Degree Overflow Handler:}\\
      Shift left producing $x^8$\\
      $\implies$ Drop $x^8$ and $\xor\ \texttt{0x1B}$.
    };
  \end{scope}
\end{tikzpicture}
\caption{The algebraic architecture of $\F_{2^8}$ (AES byte field). Bytes are interpreted directly as degree-$7$ polynomials over $\F_2$. Addition is purely carry-less bitwise XOR. Multiplication is polynomial multiplication reduced modulo the irreducible Rijndael polynomial $m(x) = x^8 + x^4 + x^3 + x + 1$, where any overflow degree $x^8$ is replaced by $x^4+x^3+x+1$ (\texttt{0x1B}).}
\label{ch06:fig-f28-poly}
\end{figure}

\subsection{Addition is free}

Adding two elements of $\F_{2^8}$ adds the polynomials coefficient-wise mod
$2$, which is bitwise XOR. There is no carry, no reduction needed (the degree
never grows). So
\[
\texttt{0x53} \xor \texttt{0xCA} = \texttt{0x99}.
\]
In hardware this is eight independent XOR gates. Subtraction is identical to
addition, since $-1 = 1$.

\begin{applicationbox}
This is exactly why $\mathrm{GF}(2^n)$ is beloved by hardware designers and
why AES is fast. Addition is XOR (no carry propagation, the slow part of
ordinary integer addition), and AES's \texttt{MixColumns} step is a small
matrix multiply over $\F_{2^8}$ that becomes a handful of XORs and table
lookups. The Galois-counter mode (GCM) used in TLS computes its
authentication tag entirely with multiplication in $\F_{2^{128}}$ --- same
construction, a degree-128 irreducible instead of degree-8 --- precisely
because field multiplication of bit-strings pipelines beautifully in silicon
(the \texttt{PCLMULQDQ} carry-less-multiply instruction exists for exactly
this).
\end{applicationbox}

\subsection{Multiplication, worked by hand}

Multiplication is where the polynomial quotient structure comes into play. The recipe:
\begin{enumerate}
  \item Multiply the two polynomials over $\F_2$ (a carry-less multiply ---
        partial products combined with XOR, not addition).
  \item If the result has degree $\geq 8$, reduce it modulo $m(x)$ by
        polynomial long division, keeping the remainder.
\end{enumerate}

Let us compute \texttt{0x57} $\cdot$ \texttt{0x83} in $\F_{2^8}$, the worked
example from the original Rijndael documentation.

\texttt{0x57} $= x^6 + x^4 + x^2 + x + 1$ and \texttt{0x83} $= x^7 + x + 1$.

\textbf{Step 1: carry-less polynomial product.} Distribute every term of the
first factor across $x^7 + x + 1$ and XOR like terms (any power that appears
an even number of times cancels mod $2$):
\[
\begin{aligned}
(x^6+x^4+x^2+x+1)\cdot x^7
  &= x^{13}+x^{11}+x^{9}+x^{8}+x^{7},\\
(x^6+x^4+x^2+x+1)\cdot x
  &= x^{7}+x^{5}+x^{3}+x^{2}+x,\\
(x^6+x^4+x^2+x+1)\cdot 1
  &= x^{6}+x^{4}+x^{2}+x+1.
\end{aligned}
\]
XOR the three rows. The two $x^7$ terms cancel, the two $x^2$ terms cancel,
and the two $x$ terms cancel, leaving
\[
x^{13}+x^{11}+x^{9}+x^{8}+x^{6}+x^{5}+x^{4}+x^{3}+1.
\]
This has degree $13$, so we must reduce.

\textbf{Step 2: reduce mod $m(x) = x^8+x^4+x^3+x+1$.} The reduction rule
comes from $m(x)=0$, i.e.
\[
x^8 = x^4 + x^3 + x + 1.
\]
Knock down each high power using $x^{k} = x^{k-8}\,(x^4+x^3+x+1)$:
\[
\begin{aligned}
x^{13} &= x^{5}(x^4{+}x^3{+}x{+}1) = x^{9}+x^{8}+x^{6}+x^{5},\\
x^{11} &= x^{3}(x^4{+}x^3{+}x{+}1) = x^{7}+x^{6}+x^{4}+x^{3},\\
x^{9}  &= x\,(x^4{+}x^3{+}x{+}1) = x^{5}+x^{4}+x^{2}+x,\\
x^{8}  &= x^4+x^3+x+1.
\end{aligned}
\]
Substitute these for the degree-$\geq 8$ terms in the dividend and XOR
everything together (the new $x^9$ and $x^8$ produced on the right then get
reduced in turn). Carrying the bookkeeping through, all powers above $7$
cancel and the surviving remainder is
\[
x^7 + x^6 + 1 \;=\; \texttt{0xC1}.
\]
So in $\F_{2^8}$,
\[
\texttt{0x57} \cdot \texttt{0x83} = \texttt{0xC1}.
\]

\begin{walkbox}
The cleanest way to do Step 2 on paper is schoolbook long division, except
``subtract'' is XOR. Write the degree-13 dividend. Its leading term is
$x^{13}$; to cancel it, XOR in $x^5 \cdot m(x)$ (which starts $x^{13}$).
That kills $x^{13}$ and flips some lower bits. Whatever is now the leading
term, cancel it with the appropriate shift $x^j \cdot m(x)$. Repeat until the
leading term has degree $< 8$; what remains is the answer. Each step is one
XOR of a shifted copy of the $9$-bit pattern \texttt{1\,0001\,1011}
($=\texttt{0x11B}$).
\end{walkbox}

\subsection{The doubling trick: \texttt{xtime}}

Full multiplication is fiddly, but multiplication by the single element
$x$ (the byte \texttt{0x02}) is so simple it has a name: \texttt{xtime}.
Multiplying a polynomial by $x$ shifts every coefficient up by one degree ---
a left shift of the byte. If the top bit $b_7$ was $0$, the result still fits
in $8$ bits and we are done. If $b_7$ was $1$, the shift produces an $x^8$
term, which we immediately reduce using $x^8 = x^4+x^3+x+1$, i.e.\ we XOR in
\texttt{0x1B}.

\[
\texttt{xtime}(b) =
\begin{cases}
b \ll 1, & \text{if top bit of } b \text{ is } 0,\\[2pt]
(b \ll 1) \xor \texttt{0x1B}, & \text{if top bit of } b \text{ is } 1.
\end{cases}
\]

\begin{example}[\texttt{xtime} in action]
\label{ch06:ex-xtime}
$\texttt{xtime}(\texttt{0x57})$: $\texttt{0x57} = 0101\,0111$, top bit $0$,
so just shift: $1010\,1110 = \texttt{0xAE}$.

$\texttt{xtime}(\texttt{0xAE})$: $\texttt{0xAE} = 1010\,1110$, top bit $1$.
Shift to $1\,0101\,1100$, drop to a byte $\texttt{0x5C}$, then XOR
\texttt{0x1B}: $\texttt{0x5C} \xor \texttt{0x1B} = \texttt{0x47}$. So
$\texttt{0x57}\cdot \texttt{0x04} = \texttt{0x47}$, since multiplying by
$\texttt{0x04} = x^2$ is two doublings.
\end{example}

Why does this matter? Because \emph{any} multiplication in $\F_{2^8}$ can be
built from \texttt{xtime} and XOR. To multiply $b$ by an arbitrary constant
$c$, write $c$ in binary; each set bit of $c$ corresponds to a power $x^k$,
and $b \cdot x^k$ is $k$ repeated \texttt{xtime}s. XOR together the
contributions for the set bits. This is ``Russian-peasant'' multiplication,
carry-less. \autoref{ch06:fig-xtime} traces the doubling chain for
$b = \texttt{0x57}$, the one the next example reuses.

\begin{figure}[H]
\centering
\begin{tikzpicture}[font=\small,>=stealth,scale=0.92]
  % 5 Chain Boxes
  \node[draw=sourceblue!80,fill=sourceblue!10,thick,rounded corners=3pt,minimum width=2.0cm,minimum height=1.3cm,align=center] (b0) at (0,0) {
    \textbf{\texttt{0x57}}\\
    \tiny $0101\,0111$\\
    \tiny $b_7 = 0$
  };
  \node[below=2pt,font=\tiny\bfseries,sourceblue] at (b0.south) {$b \cdot 1$};

  \node[draw=sourceblue!80,fill=sourceblue!10,thick,rounded corners=3pt,minimum width=2.0cm,minimum height=1.3cm,align=center] (b1) at (2.9,0) {
    \textbf{\texttt{0xAE}}\\
    \tiny $1010\,1110$\\
    \tiny \textbf{$b_7 = 1$ (overflow)}
  };
  \node[below=2pt,font=\tiny\bfseries,sourceblue] at (b1.south) {$b \cdot x$};

  \node[draw=quotientorange!80,fill=quotientorange!12,thick,rounded corners=3pt,minimum width=2.0cm,minimum height=1.3cm,align=center] (b2) at (5.8,0) {
    \textbf{\texttt{0x47}}\\
    \tiny $0100\,0111$\\
    \tiny $b_7 = 0$
  };
  \node[below=2pt,font=\tiny\bfseries,quotientorange] at (b2.south) {$b \cdot x^2$};

  \node[draw=sourceblue!80,fill=sourceblue!10,thick,rounded corners=3pt,minimum width=2.0cm,minimum height=1.3cm,align=center] (b3) at (8.7,0) {
    \textbf{\texttt{0x8E}}\\
    \tiny $1000\,1110$\\
    \tiny \textbf{$b_7 = 1$ (overflow)}
  };
  \node[below=2pt,font=\tiny\bfseries,sourceblue] at (b3.south) {$b \cdot x^3$};

  \node[draw=quotientorange!80,fill=quotientorange!12,thick,rounded corners=3pt,minimum width=2.0cm,minimum height=1.3cm,align=center] (b4) at (11.6,0) {
    \textbf{\texttt{0x07}}\\
    \tiny $0000\,0111$\\
    \tiny $b_7 = 0$
  };
  \node[below=2pt,font=\tiny\bfseries,quotientorange] at (b4.south) {$b \cdot x^4$};

  % Transition arrows
  \draw[->,very thick,sourceblue] (b0) -- node[above,font=\tiny\bfseries]{$\ll 1$} (b1);
  \draw[->,very thick,quotientorange] (b1) -- node[above,font=\tiny\bfseries,align=center]{$\ll 1$\\$\oplus\,\texttt{1B}$} (b2);
  \draw[->,very thick,sourceblue] (b2) -- node[above,font=\tiny\bfseries]{$\ll 1$} (b3);
  \draw[->,very thick,quotientorange] (b3) -- node[above,font=\tiny\bfseries,align=center]{$\ll 1$\\$\oplus\,\texttt{1B}$} (b4);
\end{tikzpicture}
\caption{The \texttt{xtime} doubling chain for $b = \texttt{0x57}$. Each step shifts the byte left by one position. Whenever the high bit $b_7$ is 1 (overflowing into $x^8$), the operation reduces modulo the Rijndael polynomial by XORing with \texttt{0x1B}. Any field multiplication reduces to selecting and XORing appropriate terms from this chain.}
\label{ch06:fig-xtime}
\end{figure}

\begin{example}[$\texttt{0x57}\cdot\texttt{0x13}$ via \texttt{xtime}]
\label{ch06:ex-mul13}
$\texttt{0x13} = x^4 + x + 1$, so we need $b\cdot x^4$, $b\cdot x$, and
$b\cdot 1$ for $b=\texttt{0x57}$, then XOR them.
\[
\begin{aligned}
b\cdot 1 &= \texttt{0x57},\\
b\cdot x &= \texttt{xtime}(\texttt{0x57}) = \texttt{0xAE},\\
b\cdot x^2 &= \texttt{xtime}(\texttt{0xAE}) = \texttt{0x47},\\
b\cdot x^3 &= \texttt{xtime}(\texttt{0x47}) = \texttt{0x8E},\\
b\cdot x^4 &= \texttt{xtime}(\texttt{0x8E}) = \texttt{0x07}.
\end{aligned}
\]
Now XOR the pieces for the set bits of \texttt{0x13} ($x^4, x^1, x^0$):
\[
\texttt{0x07} \xor \texttt{0xAE} \xor \texttt{0x57} = \texttt{0xFE}.
\]
So $\texttt{0x57}\cdot\texttt{0x13} = \texttt{0xFE}$. The small constants
$\texttt{0x01},\texttt{0x02},\texttt{0x03}$ are exactly the ones AES's
\texttt{MixColumns} multiplies by, which is why the operation is so cheap.
\end{example}

\begin{implementbox}
A complete $\F_{2^8}$ multiplier in a dozen lines. \textbf{For learning
only} --- this is a textbook reference, not a production routine. The branch
on the high bit makes it variable-time, which leaks the operands through
timing; real AES implementations use constant-time tables or hardware
\texttt{AES-NI}/\texttt{PCLMULQDQ} instructions.
\begin{lstlisting}[language=Python]
def gf_mul(a, b):
    """Multiply two bytes in GF(2^8) with the AES polynomial 0x11B."""
    p = 0
    for _ in range(8):
        if b & 1:          # current low bit of b set -> add a
            p ^= a         # addition is XOR
        b >>= 1
        hi = a & 0x80      # remember top bit before shifting
        a = (a << 1) & 0xFF
        if hi:             # overflow past x^7 -> reduce mod 0x11B
            a ^= 0x1B
    return p

def xtime(b):
    """Multiply by x (the element 0x02)."""
    b <<= 1
    if b & 0x100:
        b ^= 0x11B
    return b & 0xFF

assert gf_mul(0x57, 0x83) == 0xC1   # the worked example above
assert gf_mul(0x57, 0x13) == 0xFE
assert xtime(0x57) == 0xAE
\end{lstlisting}
The outer loop walks the bits of $b$; \texttt{xtime}'s reduction (XOR
\texttt{0x1B}) happens inline each time \texttt{a} would overflow past degree
$7$. The asserts reproduce the hand computations, so you can check the paper
work against the code.
\end{implementbox}

\subsection{Finding inverses in $\F_{2^8}$}

AES's \texttt{SubBytes} step (the S-box) is built from the multiplicative
inverse in $\F_{2^8}$: each nonzero byte is replaced by its inverse (and $0$
maps to $0$), followed by an affine bit transformation. So we need to invert
field elements. Two routes, mirroring the $\F_p$ case:

\begin{itemize}
  \item \textbf{Extended Euclid on polynomials.} Run the extended Euclidean
        algorithm on the byte's polynomial and $m(x)$; the cofactor is the
        inverse. This is the polynomial analogue of
        \autoref{ch06:thm-fp}'s proof.
  \item \textbf{Fermat-style exponentiation.} The nonzero elements form a
        group of order $255$, so $a^{255} = 1$ for every nonzero $a$, hence
        $a^{-1} = a^{254}$. Compute $a^{254}$ by repeated squaring and
        multiplying in $\F_{2^8}$.
\end{itemize}

\begin{example}[An inverse]
\label{ch06:ex-inverse}
In $\F_{2^8}$ with the AES polynomial, $\texttt{0x53}^{-1} = \texttt{0xCA}$.
You can confirm it by multiplying: running the long-division recipe on
$\texttt{0x53}\cdot\texttt{0xCA}$ collapses to the single byte $1$. Try it
once by hand and you will appreciate why implementations precompute the
answer.
\end{example}

\begin{trapbox}
Do not eyeball field inverses, and do not confuse a \emph{factor} in a
product with an \emph{inverse}. The only reliable way to get $a^{-1}$ in
$\mathrm{GF}(2^8)$ is to run extended Euclid or compute $a^{254}$ --- or to
look it up in the precomputed $256$-byte inverse table that every real AES
implementation stores. The S-box you see in the AES spec is precisely that
table (inverse, then an affine map) baked in, so the field arithmetic happens
once, offline, and the cipher itself just indexes.
\end{trapbox}

\begin{implementbox}
Compute the inverse via $a^{254}$. \textbf{Learning only.}
\begin{lstlisting}[language=Python]
def gf_inv(a):
    """Multiplicative inverse in GF(2^8): a^254, since |group| = 255."""
    if a == 0:
        return 0          # AES defines inv(0) = 0 by convention
    result, base, e = 1, a, 254
    while e:              # square-and-multiply for exponent 254
        if e & 1:
            result = gf_mul(result, base)
        base = gf_mul(base, base)
        e >>= 1
    return result

assert gf_inv(0x53) == 0xCA           # 0x53 * 0xCA == 1 in GF(2^8)
assert gf_mul(0x53, gf_inv(0x53)) == 1
\end{lstlisting}
The asserts pin down the actual inverse value, which is the only trustworthy
way to know it.
\end{implementbox}

\section{The Multiplicative Group Is Cyclic}
\label{ch06:cyclic}

The nonzero elements of any finite field form a group under multiplication,
written $\F_q^{*}$, of size $q-1$. The deepest structural fact about finite
fields, and the one cryptography exploits constantly, is this.

\begin{theorem}[The multiplicative group is cyclic]
\label{ch06:thm-cyclic}
For any finite field $\F_q$, the group $\F_q^{*}$ is cyclic of order $q-1$.
That is, there exists a \emph{generator} (or \emph{primitive element})
$g$ such that every nonzero element is a power of $g$:
\[
\F_q^{*} = \set{\, g^0, g^1, g^2, \dots, g^{q-2} \,}.
\]
\end{theorem}

\begin{intuitionbox}
``Cyclic'' means one element, raised to successive powers, cycles through
\emph{all} the nonzero elements before returning to $1$. It is the
multiplicative analogue of a clock: a single generator's powers tick through
every position. This is what makes discrete logarithms meaningful --- if the
group were a tangled product of small cycles, ``the exponent that produces
$h$'' might not be well defined or might be easy to find.
\end{intuitionbox}

\begin{example}[A generator of $\F_7^{*}$]
\label{ch06:ex-gen7}
$\F_7^{*} = \set{1,2,3,4,5,6}$, order $6$. Try $g=3$:
\[
3^1=3,\ 3^2=2,\ 3^3=6,\ 3^4=4,\ 3^5=5,\ 3^6=1.
\]
The powers of $3$ run through all six nonzero elements, so $3$ is a
generator. By contrast $2$ is not: $2^1=2, 2^2=4, 2^3=1$, then it repeats ---
$2$ has order $3$, only reaching half the group. The orders of elements
always divide $q-1=6$, here landing on $1,2,3,6$.
\end{example}

\begin{example}[A generator of $\F_8^{*}$]
\label{ch06:ex-gen8}
Build $\F_8 = \F_2[x]/(x^3+x+1)$, so $x^3 = x+1$. Is $x$ a generator? Its
order must divide $7$, which is prime, so the order is $1$ or $7$; since
$x \neq 1$, the order is $7$ and $x$ generates everything:
\[
x^1=x,\quad x^2,\quad x^3=x+1,\quad x^4 = x^2+x,\quad x^5 = x^2+x+1,
\]
\[
x^6 = x^2+1,\quad x^7 = 1.
\]
All seven nonzero elements appear exactly once. Whenever $q-1$ is prime
(here $7$), \emph{every} non-identity element is automatically a generator.
\autoref{ch06:fig-cycle8} lays the seven powers out as a clock: stepping once
around multiplies by $x$, and after seven steps you are back at $1$.
\end{example}

\begin{figure}[H]
\centering
\begin{tikzpicture}[scale=1.0,
    node distance=0pt,
    elt/.style={draw=structurepurple,fill=structurepurple!8,rounded corners=2pt,
                inner sep=2.5pt,font=\footnotesize}]
  \def\R{2.5}
  % seven positions around a circle, starting at the top, going clockwise
  \foreach \k/\lab in {0/{$1$}, 1/{$x$}, 2/{$x^2$}, 3/{$x{+}1$},
                       4/{$x^2{+}x$}, 5/{$x^2{+}x{+}1$}, 6/{$x^2{+}1$}} {
    \node[elt] (n\k) at ({90 - \k*360/7}:\R) {\lab};
  }
  % arrows: multiply by x advances one step clockwise
  \foreach \k in {0,...,6} {
    \pgfmathtruncatemacro{\next}{mod(\k+1,7)}
    \draw[-{Stealth[length=2mm]},structurepurple!70,thick]
      (n\k) to[bend right=18] (n\next);
  }
  \node[font=\footnotesize\itshape,structurepurple] at (0,0) {$\times\,x$};
\end{tikzpicture}
\caption{The cyclic group $\F_8^{*}$ generated by $x$. Each arrow is
multiplication by the generator $x$; seven steps return to $1$. Every nonzero
element sits on the cycle exactly once, which is what ``cyclic'' means.}
\label{ch06:fig-cycle8}
\end{figure}

\begin{whybox}
Why is $\F_q^{*}$ always cyclic? The clean argument: in a field, the
polynomial $t^d - 1$ has at most $d$ roots (a degree-$d$ polynomial over a
field cannot have more roots than its degree --- itself a consequence of the
division algorithm). So for each divisor $d$ of $q-1$, at most $d$ elements
satisfy $t^d = 1$. A finite abelian group with this ``at most $d$ solutions
to $t^d=1$'' property must be cyclic. The step that fails for non-fields like
$\Zn 8$ is exactly the bound on roots: $t^2-1$ has four roots in $\Zn 8$
($1,3,5,7$), more than its degree, which is why $\Zns 8$ is not cyclic.
\end{whybox}

\begin{connectionbox}
The order of an element $a$ in $\F_q^{*}$ always divides $q-1$ (Lagrange's
theorem from Chapter~5). In $\F_{2^8}$, $q-1 = 255 = 3 \cdot 5 \cdot
17$, so element orders are divisors of $255$. The element $x = \texttt{0x02}$
in the AES field happens \emph{not} to be a generator (its order is $51$, not
$255$); a primitive element there is $\texttt{0x03} = x+1$, whose powers
enumerate all $255$ nonzero bytes. Chapters~18 and~20 live
inside this cyclic structure --- there, the difficulty of \emph{reversing}
$g \mapsto g^k$ is the core asymmetry.
\end{connectionbox}

\subsection{Why AES lives in $\F_{2^8}$, summarized}
\label{ch06:why-aes}

Pulling the threads together, the AES designers wanted a byte-to-byte
nonlinear map (the S-box) that resists linear and differential cryptanalysis,
plus a column-mixing step that diffuses changes across bytes. A field gives
them exactly the algebra to make these both \emph{invertible} (decryption
must undo them) and \emph{nonlinear in the bit sense}.

\begin{itemize}
  \item \textbf{Bytes are field elements.} $2^8 = 256$ matches the byte, so
        every byte is a field element and no value is wasted.
  \item \textbf{Addition is XOR.} Free in hardware, no carries.
  \item \textbf{The inverse map $a \mapsto a^{-1}$ is highly nonlinear} over
        $\F_2$, which is precisely what frustrates linear cryptanalysis ---
        yet it is directly invertible, so decryption stays clean.
  \item \textbf{\texttt{MixColumns}} is multiplication by a fixed polynomial
        in $\F_{2^8}[y]/(y^4+1)$ with coefficients \texttt{0x01},
        \texttt{0x02}, \texttt{0x03} --- the cheap constants we saw ---
        giving strong diffusion for almost no cost.
\end{itemize}

\begin{attackbox}
What if you ``simplified'' AES by replacing $\F_{2^8}$ multiplication with
ordinary multiplication mod $256$? Then $\Zn{256}$ has zero-divisors
($2^7 \cdot 2 = 0$) and most elements have no inverse, so \texttt{MixColumns}
would not be reversible --- decryption breaks outright. Worse, the inverse
map $a\mapsto a^{-1}$ would not exist for half the bytes, gutting the S-box.
And were you to pick a \emph{reducible} degree-8 polynomial instead of an
irreducible one, the quotient $\F_2[x]/(m)$ would again have zero-divisors
(the factors of $m$ multiply to $0$), and the same collapse follows. The
irreducibility of $x^8+x^4+x^3+x+1$ is not decoration; it is what guarantees
every step of AES can be undone.
\end{attackbox}

\section{The Frobenius Map and Square Roots}
\label{ch06:frobenius}

One last structural gem, useful later for square roots, point counting on
elliptic curves, and certain attacks. Over a field of characteristic $p$,
raising to the $p$-th power is a field homomorphism.

\begin{theorem}[Freshman's dream / Frobenius]
\label{ch06:thm-frobenius}
In a field of characteristic $p$, for all $a,b$,
\[
(a+b)^p = a^p + b^p.
\]
The map $\phi(a) = a^p$ is a field automorphism of $\F_{p^n}$ (the
\emph{Frobenius} map), fixing exactly the subfield $\F_p$.
\end{theorem}

\begin{whybox}
Expand $(a+b)^p$ by the binomial theorem. Every middle coefficient
$\binom{p}{k}$ for $0<k<p$ is divisible by $p$ (because $p$ is prime and
appears in the numerator $p!$ but not the denominator $k!(p-k)!$). In
characteristic $p$, those terms vanish, leaving $a^p + b^p$. The
``mistake'' students make over $\R$ becomes a theorem over $\F_p$. In
characteristic $2$ it reads $(a+b)^2 = a^2+b^2$ --- squaring is XOR-linear,
which is why squaring in $\F_{2^n}$ is just a bit-permutation (spread the bits
out, then reduce) and cheaper than a general multiply.
\end{whybox}

\begin{example}[Squaring is linear in $\F_8$]
\label{ch06:ex-square}
Using $\F_8 = \F_2[x]/(x^3+x+1)$: $(x+1)^2 = x^2 + 1$ directly, with no cross
term $2x$ (it vanishes). And $(x^2+x)^2 = x^4 + x^2$; reduce $x^4 = x\cdot x^3
= x(x+1) = x^2 + x$, giving $x^2 + x + x^2 = x$. The Frobenius map $\phi$
sends $x \mapsto x^2$, and applying it three times returns to the start
(since $\F_8$ has $p^n = 2^3$ and $\phi^n = \mathrm{id}$).
\end{example}

\begin{connectionbox}
In characteristic $2$ the squaring map $a\mapsto a^2$ is a bijection, so
every element has a \emph{unique} square root, given by
$\sqrt{a} = a^{2^{\,n-1}}$ in $\F_{2^n}$ (because
$(a^{2^{n-1}})^2 = a^{2^n} = a$, using $a^{2^n}=a$ for all field elements).
Square roots over $\F_p$ for odd $p$ are harder --- that is the Tonelli--
Shanks algorithm, which the elliptic-curve and Rabin-cryptosystem chapters
need to recover a $y$-coordinate from $y^2$.
\end{connectionbox}

\section{Summary}
\label{ch06:summary}

\begin{itemize}
  \item A \textbf{finite field} exists for exactly the prime-power sizes
        $q = p^n$, and is unique up to isomorphism
        (\autoref{ch06:thm-classification}).
  \item $\F_p = \Zn{p}$ is a field because every nonzero residue is coprime
        to $p$ and so invertible (\autoref{ch06:thm-fp}).
  \item $\F_{p^n}$ is built as $\F_p[x]/(m(x))$ for an \textbf{irreducible}
        polynomial $m$ of degree $n$ --- the polynomial analogue of working
        modulo a prime. Irreducibility is what forces a field rather than a
        ring with zero-divisors.
  \item In $\F_{2^n}$, elements are bit-strings, addition is XOR, and
        multiplication is carry-less polynomial multiply followed by
        reduction mod $m$. \textbf{AES uses} $\F_{2^8}$ with $m(x) =
        x^8+x^4+x^3+x+1$; multiplication reduces to \texttt{xtime} (shift,
        conditional XOR \texttt{0x1B}) plus XOR.
  \item The \textbf{multiplicative group} $\F_q^{*}$ is cyclic of order
        $q-1$ (\autoref{ch06:thm-cyclic}); its generators drive discrete-log
        cryptography, and the existence of inverses is what lets AES decrypt.
  \item The \textbf{Frobenius} map $a \mapsto a^p$ is an automorphism;
        squaring in characteristic $2$ is linear, cheap, and gives unique
        square roots.
\end{itemize}

With finite fields in hand, every later building block --- AES, GCM,
Diffie--Hellman, elliptic curves, secret sharing --- has the arithmetic floor
it stands on.

Now that our algebraic foundations are complete, we need one more mathematical language before we can define modern security: probability. In Chapter~7, we turn to probability and information theory: randomness, entropy, and the birthday bound, laying the groundwork for Shannon's definition of perfect secrecy and the symmetric ciphers of Part~II.

\section*{Exercises}
\addcontentsline{toc}{section}{Exercises}

\begin{exercise}[Inverses in $\F_{11}$]
\label{ch06:ex-inv11}
Compute the multiplicative inverse of each of $2,3,5,7$ in $\F_{11}$ using
the extended Euclidean algorithm. Verify each by multiplying back to $1$.
Which element of $\F_{11}^{*}$, besides $1$, is its own inverse?
\end{exercise}

\begin{exercise}[No field of size 6]
\label{ch06:ex-no6}
Prove directly that there is no field with exactly $6$ elements. (Hint: what
must the characteristic be, and what does
\autoref{ch06:thm-classification}'s ``why $p^n$'' argument force?)
\end{exercise}

\begin{exercise}[Build $\F_9$]
\label{ch06:ex-f9}
The polynomial $x^2+1$ is irreducible over $\F_3$ (it has no root in
$\set{0,1,2}$ --- check). Construct $\F_9 = \F_3[x]/(x^2+1)$, list its $9$
elements, and compute $(x+1)(x+2)$ and $(2x+1)^{-1}$ inside it.
\end{exercise}

\begin{exercise}[Irreducible hunting]
\label{ch06:ex-irred}
Find all monic irreducible polynomials of degree $4$ over $\F_2$. There are
$16$ monic degree-4 polynomials; remove the ones with a root in $\F_2$ and
the one that is the square of the irreducible quadratic $x^2+x+1$. How many
irreducibles remain, and which are they?
\end{exercise}

\begin{exercise}[Multiply in the AES field]
\label{ch06:ex-aesmul}
Compute the following products in $\F_{2^8}$ with $m(x)=x^8+x^4+x^3+x+1$,
showing the carry-less product and the reduction:
(a) $\texttt{0x02}\cdot\texttt{0x87}$;\quad
(b) $\texttt{0x03}\cdot\texttt{0x6E}$;\quad
(c) $\texttt{0x57}\cdot\texttt{0x13}$ (you should get \texttt{0xFE}).
Check each against the \texttt{gf\_mul} routine.
\end{exercise}

\begin{exercise}[\texttt{xtime} chain]
\label{ch06:ex-xtimechain}
Starting from \texttt{0x01}, apply \texttt{xtime} repeatedly. How many steps
until you return to \texttt{0x01}? What does that number tell you about the
order of the element $x = \texttt{0x02}$ in $\F_{2^8}^{*}$, and is $x$ a
generator? (The chapter states its order; confirm it.)
\end{exercise}

\begin{exercise}[Generator search]
\label{ch06:ex-gensearch}
In $\F_{2^8}$, verify that $\texttt{0x03} = x+1$ is a primitive element by
checking that its order is $255$ and not a proper divisor of $255$. It
suffices to check $\texttt{0x03}^{255/r}\neq 1$ for $r\in\set{3,5,17}$ ---
explain why those three checks are enough.
\end{exercise}

\begin{exercise}[Break a broken cipher]
\label{ch06:ex-break}
Suppose a student implements ``AES'' but replaces the irreducible
$x^8+x^4+x^3+x+1$ with the \emph{reducible} polynomial $x^8+1 = (x+1)^8$ over
$\F_2$. Exhibit two nonzero bytes whose product is $0$ in this broken
``field,'' and explain in one sentence why \texttt{MixColumns} can no longer
be inverted.
\end{exercise}

\begin{exercise}[Frobenius cycles]
\label{ch06:ex-frob}
In $\F_{16} = \F_2[x]/(x^4+x+1)$, compute the Frobenius orbit of $x$ under
$\phi(a)=a^2$: that is, $x, x^2, x^4, x^8, x^{16}$. Confirm $x^{16}=x$ and
that the orbit has length $4 = \deg m$. What do the orbit lengths tell you
about the subfields of $\F_{16}$?
\end{exercise}

\begin{exercise}[Square roots via Frobenius]
\label{ch06:ex-sqrt}
Using $\sqrt{a} = a^{2^{\,n-1}}$ in $\F_{2^n}$, find the square root of
\texttt{0x53} in the AES field $\F_{2^8}$ (so compute $\texttt{0x53}^{128}$),
and verify by squaring your answer. Why does this exponent formula give the
square root?
\end{exercise}

\begin{exercise}[Counting irreducibles]
\label{ch06:ex-count}
The number of monic irreducible polynomials of degree $n$ over $\F_2$ is
$\frac{1}{n}\sum_{d \mid n}\mu(n/d)\,2^{d}$ (the necklace / M\"obius
formula). Use it to confirm there are $30$ irreducibles of degree $8$ over
$\F_2$ --- the pool the AES designers chose $m(x)$ from --- and write down
the count for degree $128$ (the field GCM uses) as a clean expression.
\end{exercise}
